\section{Methodology}
\label{sec:method}

\subsection{Synthetic Story Generation}
\label{sec:storygen}

We generate synthetic narratives with a procedural story generator that gives us precise control over the number of characters, attributes, and state changes.
Each story follows a three-phase structure:

\para{Introduction phase.}
Each of $N$ characters is introduced with a name, an initial location, and an initial possession.
For example: \emph{``Viktor was in the market, carrying a tattered map. Diana was in the workshop, carrying a crystal vial.''}

\para{Event phase.}
Each character undergoes a fixed number of state changes (2 in the main experiment, 4 in the stress test).
Events are either location changes (\emph{``Viktor went to the courtyard''}) or possession swaps (\emph{``Diana set down a crystal vial. Diana picked up a golden key''}).
Events are interleaved across characters to create a realistic narrative flow.

\para{Query phase.}
After the full story, we ask the model to report the final location and possession of every character simultaneously via structured JSON output.

\para{Name and attribute pools.}
We use 72 distinctive, cross-culturally diverse character names, 32 distinct locations, and 32 distinct possessions.
All names and attributes are chosen to be easily distinguishable, minimizing the chance that errors stem from lexical confusion.
Ground truth for every query is verified programmatically, with zero ambiguity.

\subsection{Experimental Conditions}

We run three experimental conditions, summarized in \tabref{tab:conditions}.

\begin{table}[t]
\centering
\caption{Experimental conditions. Each condition varies character count $N$, the number of state changes per character, and the number of stories per $N$ value.}
\label{tab:conditions}
\small
\begin{tabular}{@{}lccc@{}}
\toprule
\textbf{Condition} & \textbf{Character counts} & \textbf{State changes} & \textbf{Stories per $N$} \\
\midrule
Main & 2, 4, 8, 16, 32 & 2 & 5 \\
Extended & 2, 4, 8, 16, 32, 48, 64 & 2 & 3 \\
Stress & 4, 8, 16, 32 & 4 & 3 \\
\bottomrule
\end{tabular}
\end{table}

The main experiment tests 5 character counts with 5 stories each, yielding 620 total questions across 2 models.
The extended experiment pushes to 64 characters with 3 stories per count (1,044 total questions).
The stress test doubles the state-change density to examine the interaction between tracking load and update complexity.

\subsection{Models}

We evaluate two frontier models from the GPT-4.1 family: \gptfull (the flagship model) and \gptmini (a smaller, optimized variant).
All generations use temperature 0 and seed 42 for reproducibility, with max tokens scaled from 256 (for 2~characters) to 8,192 (for 64~characters).

\subsection{Baselines}

We compute two baselines to contextualize model performance:

\para{\randombaseline baseline.}
For each question, guess uniformly at random from all attribute values mentioned in the story.
This yields 5.7--9.7\% accuracy depending on the number of characters (more characters $=$ more possible values $=$ lower random accuracy).

\para{\initialstate baseline.}
Always predict each character's initial attribute (the value from the introduction phase, before any state changes).
This yields 28.4--29.4\% accuracy, since roughly one in three or four attributes happens not to change.

\subsection{Evaluation Metrics}

\para{Accuracy.}
We compute exact-match accuracy: the fraction of (character, attribute) queries for which the model's answer matches the ground truth.
We report 95\% bootstrap confidence intervals for all accuracy estimates.

\para{Error taxonomy.}
We classify every incorrect answer into one of three categories:
\begin{itemize}[leftmargin=*,itemsep=0pt,topsep=0pt]
\item \textbf{Confusion}: The model's answer is the correct current value for a \emph{different} character. This indicates cross-character interference.
\item \textbf{Forgetting}: The model's answer matches the character's \emph{initial} attribute (before state changes). This indicates failure to register updates.
\item \textbf{Other}: The answer is neither another character's current value nor the queried character's initial value (\eg returning ``none'' or a hallucinated value).
\end{itemize}

\para{Position analysis.}
We compute Spearman rank correlation between a character's introduction order (position in the story) and accuracy, to test for primacy/recency effects consistent with the ``lost in the middle'' phenomenon~\citep{liu2024lost}.
